\subsection{Projective Garbling}
\label{sec:garbling}

Bellare, Hoang, and Rogaway~\cite{BellareHoangRogaway2012Foundations}
gave a systematic formalization of garbling schemes as a standalone
cryptographic primitive and of their security notions.  We use the
three-algorithm presentation of Hemenway et al.\
\cite[full version, Section~3, Definition~1]{HemenwayEtAl2016AdaptiveGarbling}, which
fits our separation between preparing a garbled circuit and later
releasing its encoded input.  We require projectivity and adaptive
simulation privacy.

\subsubsection{Syntax and correctness}
Let $C:\bits^n\to\bits^m$ be a Boolean circuit.  A garbling scheme
consists of three polynomial-time algorithms:
\begin{description}
  \item[$\Garble(1^\lambda,C)\to(\GC,e)$.]
    Randomly generate a public garbled circuit $\GC$ and a private
    input-encoding key $e$.
  \item[$\Encode(e,x)\to X$.]
    Deterministically encode an input $x\in\bits^n$ as a garbled input $X$.
  \item[$\Eval(\GC,X)\to y$.]
    Deterministically evaluate the garbled circuit, returning
    $y\in\bits^m$ or $\bot$.
\end{description}
Correctness requires that, for every polynomial-size circuit family and
every corresponding input family
\[
  \Pr_{(\GC,e)\sample\Garble(1^\lambda,C)}
  \bigl[\Eval(\GC,\Encode(e,x))=C(x)\bigr] \ge 1-\negl(\lambda).
\]

The scheme is \emph{projective} if $e$ specifies two labels
$L_{w,0},L_{w,1}$ for each input wire $w\in\{1,\ldots,n\}$ and
possibly common auxiliary material $\GarbleAux$, such that
\[
  \Encode(e,x)
    =\bigl(\GarbleAux,L_{1,x_1},\ldots,L_{n,x_n}\bigr).
\]
Labels and auxiliary material are polynomial-length strings; we do not
require them to have length $\lambda$.  For fixed $e$, the common part
$\GarbleAux$ is independent of $x$ and may be empty.  Only $\GC$ is
released initially; $\GarbleAux$ is released with the selected labels.
It may contain output-decoding information or a key needed to unlock
the garbled circuit.  Separating this common part is a notational
convenience, not an additional security property.

\subsubsection{Adaptive privacy}
Write $\Topo(C)$ for the circuit's wiring and designated input
and output wires, without the functions computed by its gates.  We allow
this topology to be public.  A simulator consists of two efficient
algorithms $\mathcal S=(\SimGarble,\SimEncode)$.

The experiments $\GarbleReal_{\mathcal A}(1^\lambda)$ and
$\GarbleIdeal_{\mathcal A,\mathcal S}(1^\lambda)$ initially give
$1^\lambda$ to the efficient adversary $\mathcal A$ and offer these
commands:
\begin{description}
  \item[$\GarbleCircuit(C)$.]
    Store $C$.  In the real experiment, compute
    $(\GC,e)\sample\Garble(1^\lambda,C)$.
    In the ideal experiment, compute
    $(\GC,\SimState)\sample
    \SimGarble(1^\lambda,\Topo(C))$.
    Return $\GC$, retaining $e$ or $\SimState$ privately.
  \item[$\EncodeInput(x)$.]
    For $x\in\bits^n$, compute $X\gets\Encode(e,x)$ in the
    real experiment.  In the ideal experiment, compute
    $X\sample\SimEncode(\SimState,C(x))$.  Return $X$.
  \item[$\Finish(b)$.]
    Terminate and output $b\in\bits$.
\end{description}
There is at most one accepted $\GarbleCircuit$ command, followed
by at most one accepted $\EncodeInput$ command.  Malformed,
out-of-order, or repeated commands return $\bot$ without changing state.
The adversary may terminate early; termination without an output bit is
interpreted as $\Finish(0)$.

\begin{definition}[Adaptive garbling privacy]
\label{def:adaptive-garbling}
A garbling scheme has adaptive simulation privacy if there exists an
efficient simulator $\mathcal S$ such that, for every efficient adversary
$\mathcal A$, there is a negligible function $\varepsilon_{\mathcal A}$
satisfying
\[
  \left|
    \Pr[\GarbleReal_{\mathcal A}(1^\lambda)=1]
    -\Pr[\GarbleIdeal_{\mathcal A,\mathcal S}(1^\lambda)=1]
  \right|\leq\varepsilon_{\mathcal A}(\lambda).
\]
\end{definition}

Thus the simulator produces $\GC$ before learning the output, and later
produces $X$ from $C(x)$ alone, without receiving $x$ or the gate
functions.  The adversary chooses $x$ after seeing $\GC$, but chooses
all of $x$ before receiving $X$.  This is the coarse-grained adaptive
privacy notion $\PrvOne$ of Bellare, Hoang, and Rogaway~
\cite{BellareHoangRogaway2012Adaptive}; it does not permit several
input encodings for the same garbling.  Here adaptivity concerns input
selection, not corruption of a garbler.

\subsubsection{Existence}
\begin{theorem}[Hemenway et al.]
\label{thm:adaptive-garbling}
If one-way functions exist, then projective garbling schemes with
adaptive simulation privacy exist in the standard model.
\end{theorem}
This follows from the projective construction and
\cite[full version, Corollary~2]{HemenwayEtAl2016AdaptiveGarbling}.  We use the scheme
as a black box.
